\documentclass[profile=standard]{ipara-handbook}
\title{B01｜局部线性化：微分与线性映射}
\author{}
\date{}
\begin{document}
\maketitle

\begin{quote}
这座桥回答一个跨学科母问题：\textbf{高数里“可微”的真正输出为什么不是一堆偏导数，而是一张线性映射？这张线性映射又怎样被线代的矩阵语言接收？}
\end{quote}

\section{本册定位：把“变化率”升级为“统一线性主部”}
高数 Topic07 Own 多元可微、偏导、方向导数、梯度、Jacobian matrix 与链式法则；线代 Topic02 Own 线性映射及其矩阵表示。B01 不重新定义两侧，而只 Own 一条接口：
\[
\boxed{
F\text{ 在 }x\text{ 可微}
\Longrightarrow
\text{存在唯一线性映射 }DF_x
\Longrightarrow
\text{选基后得到矩阵 }J_F(x)
}
\]
这里第一步是数学蕴含，第二步是“对象到表示”的接口翻译。

若没有这座桥，最常见的断裂是：学生会算 $F_x,F_y$，却不知道为什么它们应该组织成一张矩阵；会背 Jacobian matrix，却把它误认为原非线性函数本身；会做链式法则，却说不清为什么矩阵乘法次序不能反。

\section{为什么偏导数不够：它们只是坐标轴探针}
设 $F:\mathbb R^n\to\mathbb R^m$。偏导数只观察
\[
F(x+t e_i)-F(x)
\]
沿第 $i$ 个坐标轴的变化。方向导数则把探针扩展到某条固定方向 $v$。但“每条方向都能单独测出变化率”仍不等于存在一个统一模型同时解释所有小位移。

真正的可微要求：存在一个线性映射 $L:\mathbb R^n\to\mathbb R^m$，使
\[
F(x+h)=F(x)+Lh+r(h),
\qquad
\frac{\lVert r(h)\rVert}{\lVert h\rVert}\to0.
\]
此时 $L$ 唯一，记为 $DF_x$。

\subsection{“最佳局部线性”到底最佳在哪里}
不是说 $DF_xh$ 在每个有限尺度上都最接近 $F(x+h)-F(x)$，而是说误差相对于步长是一阶更小量：
\[
F(x+h)-F(x)-DF_xh=o(\lVert h\rVert).
\]
所以当 $h$ 缩小一倍，主部随 $\lVert h\rVert$ 一阶缩小，而余项比这一尺度更快消失。这就是“线性主部支配局部变化”的严格含义。

\section{线代怎样接收：矩阵只是 $DF_x$ 的坐标记录}
在线代中，一个抽象线性映射只有在选定定义域与陪域的基以后，才写成矩阵。对标准基，$DF_x$ 的矩阵就是 Jacobian matrix：
\[
J_F(x)
=
\begin{pmatrix}
\dfrac{\partial F_1}{\partial x_1}&\cdots&\dfrac{\partial F_1}{\partial x_n}\\
\vdots&&\vdots\\
\dfrac{\partial F_m}{\partial x_1}&\cdots&\dfrac{\partial F_m}{\partial x_n}
\end{pmatrix}.
\]
因此
\[
DF_x(h)\quad\leftrightarrow\quad J_F(x)[h]_{
m std}.
\]
这一步必须记住对象边界：
\[
\boxed{DF_x\neq J_F(x),\qquad J_F(x)\text{ 是 }DF_x\text{ 在给定基下的表示}.}
\]

\section{为什么链式法则必然变成矩阵乘法}
若 $F:\mathbb R^n\to\mathbb R^m$、$G:\mathbb R^m\to\mathbb R^p$ 都可微，则局部先经过 $F$ 的线性主部，再经过 $G$ 的线性主部：
\[
D(G\circ F)_x=DG_{F(x)}\circ DF_x.
\]
选标准基后得到
\[
J_{G\circ F}(x)=J_G(F(x))J_F(x).
\]
乘法顺序不是口诀，而是函数复合顺序的坐标表达：向量 $h$ 先被 $J_F$ 作用，再被 $J_G$ 作用。

\subsection{维度是最便宜的独立校验}
若
\[
J_F:m\times n,
\qquad
J_G:p\times m,
\]
则
\[
J_GJ_F:p\times n,
\]
正好对应复合映射 $\mathbb R^n\to\mathbb R^p$。如果维度乘不上，说明对象顺序已经错了。

\section{一个母例：同一个局部线性对象怎样被不同下游继续读取}
设
\[
F(x,y)=\bigl(x^2+xy,e^y\bigr).
\]
在 $(1,0)$ 处
\[
J_F(1,0)=
\begin{pmatrix}
2&1\\0&1
\end{pmatrix}.
\]
所以
\[
F(1+h,k)=F(1,0)+
\begin{pmatrix}
2&1\\0&1
\end{pmatrix}
\binom hk+o\!\left(\sqrt{h^2+k^2}\right).
\]
这一张局部线性对象可以被多个 Bridge 继续读取：
\begin{itemize}
\item B02 取其 determinant，解释局部面积缩放；
\item 若 $F$ 是标量函数的一阶导数对象，B04 读取其 kernel/法空间结构；
\item 二阶升级后，B03 读取 Hessian 产生的二次型；
\item 复合映射继续用线性映射复合解释链式法则。
\end{itemize}
这说明 B01 不是“可微的一种解释”，而是 B02/B03/B04 的共同上游接口。

\section{最重要的资格边界}
\subsection{偏导存在不推出可微}
偏导只给坐标轴信息。如果不同方向上的误差无法由同一个线性映射统一控制，就不能过 B01。

典型反例
\[
f(x,y)=
\begin{cases}
\dfrac{x^2y}{x^4+y^2},&(x,y)\neq(0,0),\\
0,&(0,0),
\end{cases}
\]
可以出现若干方向导数存在却不连续的现象。判断可微资格时，最终仍要回到统一余项，而不是“导数数量够不够多”。

\subsection{$C^1$ 是方便的充分条件，不是可微的定义}
偏导在邻域连续可以推出可微，但可微本身并不要求偏导函数在邻域连续。不能把“常用判据”升级成定义。

\subsection{Jacobian determinant 只在方阵情形存在}
$F:\mathbb R^n\to\mathbb R^m$ 的 Jacobian matrix 总可以是矩形。只有 $m=n$ 时，才有 ordinary determinant。B02 只接收满足该维度资格的局部线性映射。

\section{跨章节地图：B01 串起了哪些内容}
\begin{tabular}{p{0.22\linewidth}p{0.32\linewidth}p{0.38\linewidth}}
\hline
章节 & 交给 B01 的对象 & B01 输出到哪里\\
\hline
高数 Topic07 可微 & 统一小位移模型 & 线代 Topic02 的线性映射/矩阵表示\\
高数 Topic07 链式 & 两层局部变化 & 线性映射复合与矩阵乘法\\
高数 Topic08 换元 & 坐标变换的一阶作用 & B02 determinant 体积缩放\\
高数 Topic07 二阶 Taylor & 一阶模型完成后的升级入口 & B03 Hessian 二次型\\
高数 Topic07 约束 & 约束函数的一阶线性条件 & B04 kernel/法空间\\
\hline
\end{tabular}

\section{调用信号与退出边界}
题面出现“小增量近似、全微分、误差传播、复合映射、Jacobian matrix、局部坐标变化”时，先问：
\begin{enumerate}
\item 当前对象是否真的可微，而不只是偏导存在？
\item 输入空间、输出空间的维度分别是多少？
\item 线性主部 $DF_x$ 作用在哪个增量 $h$ 上？
\item 当前目标只是求导，还是要继续读 determinant、二阶形状或约束空间？
\end{enumerate}
若目标只是得到 Jacobian matrix，到这里即可退出 B01；若问体积、极值、约束，则分别转 B02/B03/B04。

\section{反桥梁：几个看起来像但不能混的对象}
\begin{tabular}{p{0.2\linewidth}p{0.04\linewidth}p{0.2\linewidth}p{0.48\linewidth}}
\hline
A & & B & 真正区别\\
\hline
偏导存在 & $\neq$ & 可微 & 前者只看坐标轴；后者要求统一线性余项\\
方向导数存在 & $\neq$ & 可微 & 每条方向分别存在，不代表方向到导数的关系线性且余项一致\\
Jacobian matrix & $\neq$ & 原映射 & 前者是点处导数的矩阵表示；后者通常非线性\\
Jacobian matrix & $\neq$ & determinant & 一个记录完整线性作用，一个只压缩方阵的体积缩放\\
\hline
\end{tabular}

\section{一页压缩}
忘掉所有坐标公式，只保留：
\[
\boxed{F(x+h)=F(x)+DF_x(h)+o(\lVert h\rVert)}.
\]
然后问四件事：
\begin{enumerate}
\item $DF_x$ 是否真的存在？
\item 它是从哪个输入空间到哪个输出空间的线性映射？
\item 当前矩阵只是怎样的坐标表示？
\item 下游究竟要读它的完整方向作用、determinant、二阶升级，还是约束 kernel？
\end{enumerate}
能从这四问重新生成 Jacobian、链式与三个下游 Bridge，B01 才真正掌握。

\end{document}
